% ========================================
% 代码高亮配置模板 (增强版)
% 适用于多种编程语言的教材
% 版本: v2.0
% 更新: 2025年8月7日
% ========================================

\usepackage{listings}
\usepackage{xcolor}

% ========================================
% 颜色定义
% ========================================

% 代码配色方案
\definecolor{codeblue}{RGB}{40, 90, 200}
\definecolor{codegreen}{RGB}{0, 150, 0}
\definecolor{codered}{RGB}{200, 0, 0}
\definecolor{codepurple}{RGB}{150, 0, 150}
\definecolor{codegray}{RGB}{128, 128, 128}
\definecolor{codeorange}{RGB}{255, 140, 0}
\definecolor{backcolour}{RGB}{248, 248, 248}
\definecolor{framecolor}{RGB}{220, 220, 220}

% ========================================
% 全局 listings 配置
% ========================================

\lstset{
    backgroundcolor=\color{backcolour},   % 背景色
    commentstyle=\color{codegreen}\itshape, % 注释样式
    keywordstyle=\color{codeblue}\bfseries, % 关键字样式
    numberstyle=\tiny\color{codegray},    % 行号样式
    stringstyle=\color{codered},          % 字符串样式
    basicstyle=\ttfamily\footnotesize,    % 基本字体
    breakatwhitespace=false,              % 只在空白处换行
    breaklines=true,                      % 自动换行
    captionpos=b,                         % 标题位置
    keepspaces=true,                      % 保持空格
    numbers=left,                         % 行号位置
    numbersep=8pt,                        % 行号与代码间距
    showspaces=false,                     % 不显示空格
    showstringspaces=false,               % 不显示字符串中的空格
    showtabs=false,                       % 不显示tab
    tabsize=4,                            % tab大小
    frame=single,                         % 边框样式
    rulecolor=\color{framecolor},         % 边框颜色
    framexleftmargin=8pt,                 % 左边距
    framexrightmargin=8pt,                % 右边距
    framextopmargin=8pt,                  % 上边距
    framexbottommargin=8pt,               % 下边距
    xleftmargin=15pt,                     % 整体左边距
    xrightmargin=15pt,                    % 整体右边距
    escapeinside={\%*}{*\)},              % 转义字符
    extendedchars=true,                   % 支持扩展字符
    literate={·}{{\textperiodcentered}}1  % 中文标点支持
}

% ========================================
% C# 语言配置
% ========================================

\lstdefinestyle{csharp}{
    language=[Sharp]C,
    morekeywords={
        var, async, await, yield, partial, get, set, value,
        public, private, protected, internal, static, virtual,
        override, abstract, sealed, readonly, const, extern,
        unsafe, volatile, implicit, explicit, operator,
        namespace, using, class, struct, interface, enum,
        delegate, event, this, base, new, typeof, sizeof,
        is, as, null, true, false, if, else, switch, case,
        default, for, foreach, while, do, break, continue,
        return, throw, try, catch, finally, checked, unchecked,
        lock, in, out, ref, params, stackalloc
    },
    morecomment=[l]{//},
    morecomment=[s]{/*}{*/},
    morecomment=[s]{///}{},
    morestring=[b]",
    morestring=[b]',
    showstringspaces=false,
    keywordstyle=\color{codeblue}\bfseries,
    commentstyle=\color{codegreen}\itshape,
    stringstyle=\color{codered},
    identifierstyle=\color{black},
    emphstyle=\color{codepurple}\bfseries,
    emph={Console, WriteLine, ReadLine, String, Int32, Double, Bool, List, Dictionary, Array},
    title=\lstname
}

% ========================================
% JavaScript/TypeScript 配置
% ========================================

\lstdefinestyle{javascript}{
    language=JavaScript,
    morekeywords={
        let, const, var, function, return, if, else, for, while,
        do, break, continue, switch, case, default, try, catch,
        finally, throw, new, this, typeof, instanceof, in, of,
        class, extends, super, static, get, set, async, await,
        import, export, from, default, as, yield, delete,
        undefined, null, true, false, NaN, Infinity
    },
    morecomment=[l]{//},
    morecomment=[s]{/*}{*/},
    morestring=[b]",
    morestring=[b]',
    morestring=[b]`,
    showstringspaces=false,
    keywordstyle=\color{codeblue}\bfseries,
    commentstyle=\color{codegreen}\itshape,
    stringstyle=\color{codered},
    identifierstyle=\color{black},
    emphstyle=\color{codepurple}\bfseries,
    emph={console, log, document, window, addEventListener, querySelector, Promise, Array, Object, JSON},
    title=\lstname
}

% ========================================
% Python 配置
% ========================================

\lstdefinestyle{python}{
    language=Python,
    morekeywords={
        and, as, assert, break, class, continue, def, del, elif,
        else, except, exec, finally, for, from, global, if,
        import, in, is, lambda, not, or, pass, print, raise,
        return, try, while, with, yield, True, False, None,
        async, await, nonlocal
    },
    morecomment=[l]{\#},
    morestring=[b]",
    morestring=[b]',
    morestring=[s]{"""}{"""},
    morestring=[s]{'''}{'''},
    showstringspaces=false,
    keywordstyle=\color{codeblue}\bfseries,
    commentstyle=\color{codegreen}\itshape,
    stringstyle=\color{codered},
    identifierstyle=\color{black},
    emphstyle=\color{codepurple}\bfseries,
    emph={self, cls, len, str, int, float, list, dict, tuple, set, range, enumerate, zip},
    title=\lstname
}

% ========================================
% Java 配置
% ========================================

\lstdefinestyle{java}{
    language=Java,
    morekeywords={
        abstract, assert, boolean, break, byte, case, catch,
        char, class, const, continue, default, do, double,
        else, enum, extends, final, finally, float, for,
        goto, if, implements, import, instanceof, int,
        interface, long, native, new, package, private,
        protected, public, return, short, static, strictfp,
        super, switch, synchronized, this, throw, throws,
        transient, try, void, volatile, while, true, false, null
    },
    morecomment=[l]{//},
    morecomment=[s]{/*}{*/},
    morestring=[b]",
    morestring=[b]',
    showstringspaces=false,
    keywordstyle=\color{codeblue}\bfseries,
    commentstyle=\color{codegreen}\itshape,
    stringstyle=\color{codered},
    identifierstyle=\color{black},
    emphstyle=\color{codepurple}\bfseries,
    emph={System, out, println, String, Integer, Double, Boolean, ArrayList, HashMap},
    title=\lstname
}

% ========================================
% SQL 配置
% ========================================

\lstdefinestyle{sql}{
    language=SQL,
    morekeywords={
        SELECT, FROM, WHERE, INSERT, UPDATE, DELETE, CREATE,
        DROP, ALTER, TABLE, DATABASE, INDEX, VIEW, PROCEDURE,
        FUNCTION, TRIGGER, JOIN, INNER, LEFT, RIGHT, OUTER,
        ON, GROUP, BY, ORDER, HAVING, UNION, DISTINCT, AS,
        AND, OR, NOT, IN, EXISTS, BETWEEN, LIKE, IS, NULL,
        TRUE, FALSE, PRIMARY, FOREIGN, KEY, CONSTRAINT,
        REFERENCES, AUTO_INCREMENT, DEFAULT, CHECK
    },
    morecomment=[l]{--},
    morecomment=[s]{/*}{*/},
    morestring=[b]',
    morestring=[b]",
    showstringspaces=false,
    keywordstyle=\color{codeblue}\bfseries,
    commentstyle=\color{codegreen}\itshape,
    stringstyle=\color{codered},
    identifierstyle=\color{black},
    basicstyle=\ttfamily\footnotesize,
    title=\lstname
}

% ========================================
% XML/HTML 配置
% ========================================

\lstdefinestyle{xml}{
    language=XML,
    morestring=[b]",
    morestring=[b]',
    morecomment=[s]{<?}{?>},
    morecomment=[s]{<!--}{-->},
    commentstyle=\color{codegreen}\itshape,
    stringstyle=\color{codered},
    identifierstyle=\color{black},
    keywordstyle=\color{codeblue}\bfseries,
    showstringspaces=false,
    basicstyle=\ttfamily\footnotesize,
    title=\lstname
}

% ========================================
% CSS 配置
% ========================================

\lstdefinestyle{css}{
    language=CSS,
    morestring=[b]",
    morestring=[b]',
    morecomment=[s]{/*}{*/},
    morekeywords={
        color, background, font, margin, padding, border,
        width, height, display, position, float, clear,
        text-align, font-size, font-family, font-weight,
        line-height, vertical-align, overflow, visibility,
        z-index, top, right, bottom, left, absolute, relative,
        fixed, static, block, inline, none, auto, inherit
    },
    keywordstyle=\color{codeblue}\bfseries,
    commentstyle=\color{codegreen}\itshape,
    stringstyle=\color{codered},
    identifierstyle=\color{black},
    showstringspaces=false,
    basicstyle=\ttfamily\footnotesize,
    title=\lstname
}

% ========================================
% 配置文件样式 (JSON, YAML, INI)
% ========================================

\lstdefinestyle{json}{
    basicstyle=\ttfamily\footnotesize,
    showstringspaces=false,
    morestring=[b]",
    stringstyle=\color{codered},
    keywordstyle=\color{codeblue}\bfseries,
    numberstyle=\color{codegray},
    commentstyle=\color{codegreen}\itshape,
    frame=single,
    rulecolor=\color{framecolor},
    backgroundcolor=\color{backcolour},
    title=\lstname
}

% ========================================
% 命令行/终端样式
% ========================================

\lstdefinestyle{terminal}{
    basicstyle=\ttfamily\footnotesize\color{white},
    backgroundcolor=\color{black},
    showstringspaces=false,
    numbers=none,
    frame=single,
    rulecolor=\color{codegray},
    xleftmargin=15pt,
    xrightmargin=15pt,
    title=\lstname
}

% ========================================
% 伪代码样式
% ========================================

\lstdefinestyle{pseudocode}{
    basicstyle=\ttfamily\footnotesize,
    showstringspaces=false,
    keywordstyle=\color{codeblue}\bfseries,
    commentstyle=\color{codegreen}\itshape,
    morekeywords={
        algorithm, begin, end, if, then, else, endif, while,
        do, endwhile, for, to, endfor, repeat, until, return,
        input, output, function, procedure, call
    },
    morecomment=[l]{//},
    frame=single,
    rulecolor=\color{framecolor},
    backgroundcolor=\color{backcolour},
    title=\lstname
}

% ========================================
% 便捷命令定义
% ========================================

% 快速代码环境
\newcommand{\cscode}[1]{\lstinline[style=csharp]!#1!}
\newcommand{\jscode}[1]{\lstinline[style=javascript]!#1!}
\newcommand{\pycode}[1]{\lstinline[style=python]!#1!}
\newcommand{\javacode}[1]{\lstinline[style=java]!#1!}
\newcommand{\sqlcode}[1]{\lstinline[style=sql]!#1!}

% 代码块环境
\lstnewenvironment{csharpcode}[1][]
{\lstset{style=csharp, #1}}
{}

\lstnewenvironment{javascriptcode}[1][]
{\lstset{style=javascript, #1}}
{}

\lstnewenvironment{pythoncode}[1][]
{\lstset{style=python, #1}}
{}

\lstnewenvironment{javacode}[1][]
{\lstset{style=java, #1}}
{}

\lstnewenvironment{sqlcode}[1][]
{\lstset{style=sql, #1}}
{}

\lstnewenvironment{xmlcode}[1][]
{\lstset{style=xml, #1}}
{}

\lstnewenvironment{terminalcode}[1][]
{\lstset{style=terminal, #1}}
{}

% ========================================
% 使用说明
% ========================================

% 在文档中使用代码的方法：
% 
% 1. 行内代码：
%    \cscode{var x = 10;}
%    \jscode{let x = 10;}
%    \pycode{x = 10}
%
% 2. 代码块：
%    \begin{csharpcode}[caption=C\# 示例]
%    using System;
%    class Program {
%        static void Main() {
%            Console.WriteLine("Hello World!");
%        }
%    }
%    \end{csharpcode}
%
% 3. 直接使用 lstlisting：
%    \begin{lstlisting}[style=python, caption=Python示例]
%    def hello():
%        print("Hello World!")
%    \end{lstlisting}
%
% 4. 从文件导入代码：
%    \lstinputlisting[style=java, caption=Java示例]{example.java}
